\documentclass[11pt]{article}
\usepackage[T1]{fontenc}
\usepackage{lmodern}
\usepackage[margin=1in]{geometry}
\usepackage{amsmath,amssymb,amsthm,mathtools,microtype}
\usepackage{booktabs}
\usepackage[hidelinks]{hyperref}
\hypersetup{
 pdftitle={Extremizers of the expected missing mass under an entropy constraint},
 pdfauthor={Christopher D. Long},
 pdfsubject={Two-size structure, finite classification, and sharp asymptotics}}
\numberwithin{equation}{section}
\allowdisplaybreaks[1]
\newtheorem{theorem}{Theorem}[section]
\newtheorem{lemma}[theorem]{Lemma}
\newtheorem{corollary}[theorem]{Corollary}
\newcommand{\dd}{\,\mathrm d}
\newcommand{\EE}{\mathbb E}
\newcommand{\K}{\mathcal K}
\newcommand{\HH}{\mathcal H}
\title{Extremizers of the expected missing mass\\under an entropy constraint}
\author{Christopher D. Long\\[3pt]
 \small\href{mailto:galizur@gmail.com}{\texttt{galizur@gmail.com}}}
\date{}
\begin{document}
\maketitle
\begin{abstract}
For a discrete probability distribution $p=(p_i)$, the expected mass of
symbols absent from $t$ independent samples is
$\Phi_t(p)=\sum_i p_i(1-p_i)^t$. For every positive integer $t$ and finite
entropy budget $h\ge0$, we prove that every $\ell^1$-local maximizer under
$H(p)\le h$ has at most two distinct positive coordinate values, and all
but at most one positive coordinate are equal. This holds on finite and
countably infinite alphabets and improves the four-size bound of Berend,
Kontorovich, and Zagdanski (2017). On a countably infinite alphabet, we give
an exact finite classification of all global maximizers: a unique
one-light or uniform candidate is compared with an explicitly bounded
family of one-heavy candidates, whose masses are uniquely specified by
entropy equations. The one-light or uniform candidate is uniquely optimal
for $t=1,2$. For fixed $h>0$, the optimal value $B_t(h)$ satisfies
\[
 \frac{h}{B_t(h)}=\log t+\log\log t+2
       +O_h\!\left(\frac{\log\log t}{\log t}\right).
\]
Every optimizer is eventually one-heavy, with $m_t\sim ht$ equal light
atoms of size $q_t\sim1/(t\log t)$. The two-size proof combines a
three-coordinate entropy-curvature estimate with a finite coefficient
comparison. Crossing theorems describe the winning intervals of the
candidate distributions, and the classification gives exact extrema for
singleton counts, discovery increments, and expected coverage.
\end{abstract}

\section{Introduction}
Let $I$ be a nonempty finite or countably infinite alphabet, and write
\[
 H(p)=-\sum_{i\in I}p_i\log p_i,\qquad
 \Phi_t(p)=\sum_{i\in I}p_i(1-p_i)^t,
\]
\[
 \mathcal P_h(I)=\left\{p_i\ge0:\ \sum_{i\in I}p_i=1,\ H(p)\le h\right\}.
\]
Here $0\log0=0$ and logarithms are natural. For integer $t\ge1$, $\Phi_t$
is the expected mass of symbols absent from $t$ independent samples.
An \emph{atom size} is a distinct positive coordinate value.

Without an entropy constraint, Berend and Kontorovich~\cite{BK} showed that
finite-alphabet local maximizers have all but possibly one coordinate equal,
with the remaining coordinate no smaller. For the entropy-constrained
problem, Berend, Kontorovich, and Zagdanski~\cite[Theorem~1]{BKZ} obtained a
four-size bound for finite-alphabet local maximizers and asked whether four
sizes can occur. They also characterized the maximizing shape for
$t=1$~\cite[Proposition~3]{BKZ}. The following theorem replaces four by two.

\begin{theorem}\label{thm:main}
For every integer $t\ge1$ and $0\le h<\infty$, every local maximizer of
$\Phi_t$ on $\mathcal P_h(I)$, in the $\ell^1$ topology, has at most two
atom sizes and has finite support. The maximum is attained. If $I$ is
countably infinite, every local maximizer satisfies $H(p)=h$.
\end{theorem}

On countably infinite alphabets, attainment and finite support of global
maximizers were established in~\cite[Theorem~2]{BKZ}, and entropy saturation
in~\cite[Proposition~1 and Corollary~1]{BKZ}. The finite-subvector argument in
their proof of Theorem~2 also yields a four-size bound, hence finite support,
for $\ell^1$-local maximizers: fixing the complement makes each finite
subvector locally optimal with its own mass and entropy budget, to which
their four-size argument applies. We include the compactness and
atom-splitting arguments for completeness.

Theorem~\ref{thm:exceptional} strengthens the two-size conclusion: every
nonuniform local maximizer has one exceptional atom and an equimass block.
For a countably infinite alphabet, Theorem~\ref{thm:finite-classification}
reduces the global problem to a finite candidate list with an explicit
multiplicity cutoff. Theorem~\ref{thm:asymptotics} determines
the next terms beyond the first-order asymptotic $h/\log t$ from
\cite[Theorems~3--4]{BKZ}, and also determines the asymptotic sizes and
multiplicities of the maximizing atoms.

The proof of the two-size theorem occupies the next three sections.
Section~\ref{sec:classification} gives the multiplicity restriction and
finite classification, and Section~\ref{sec:asymptotics} gives the
asymptotics and occupancy consequences. The appendices contain the
entropy-curvature estimate, crossing comparisons, the finite-alphabet
classification, and rationally certified examples.

\section{The three-coordinate reduction}\label{sec:three-coordinate}

Fix three distinct positive coordinates $x<y<z$, keeping all other
coordinates unchanged, and put
\begin{equation}\label{eq:geometry}
\begin{gathered}
 S=x+y+z,\quad E=xy+xz+yz,\quad P=xyz,\quad D(s)=(s+x)(s+y)(s+z),\\
 I_0=\int_0^\infty\frac{\dd s}{D(s)},\qquad
 I_1=\int_0^\infty\frac{s\dd s}{D(s)},\qquad m=I_1/I_0,\\
 \K=\frac1{I_0}\int_0^\infty\frac{(s-m)^2}{D(s)^2}\dd s.
\end{gathered}
\end{equation}
The map $(x,y,z)\mapsto(S,E,P)$ has nonzero Vandermonde Jacobian at a
triple of distinct roots. Thus, with $S$ fixed, $(E,P)$ are smooth local
coordinates; sufficiently small changes preserve positivity and distinctness.
The three-coordinate entropy $\HH=-x\log x-y\log y-z\log z$ satisfies
\begin{equation}\label{eq:entropy-derivatives}
 \HH_E=I_1,\qquad \HH_P=I_0.
\end{equation}
Indeed the roots of $W(u)=u^3-Su^2+Eu-P$ satisfy
$u_E=-u/W'(u)$ and $u_P=1/W'(u)$. Use
\[
 \log u=\int_0^\infty\left(\frac1{s+1}-\frac1{s+u}\right)\dd s
\]
and the partial fractions, with each sum over $u\in\{x,y,z\}$,
\[
 \sum_u\frac1{(s+u)W'(u)}=\frac1{D(s)},\qquad
 \sum_u\frac{u}{(s+u)W'(u)}=-\frac{s}{D(s)}.
\]
Differentiation under the integral gives
\[
 \HH_{EE}=-\int_0^\infty\frac{s^2}{D(s)^2}\dd s,\quad
 \HH_{EP}=-\int_0^\infty\frac{s}{D(s)^2}\dd s,\quad
 \HH_{PP}=-\int_0^\infty\frac{\dd s}{D(s)^2}.
\]
Since $\HH_P=I_0>0$, the entropy level set through the triple is a smooth
curve $P=P(E)$, and implicit differentiation yields
\begin{equation}\label{eq:curve}
 P'=-m,\qquad P''=\K.
\end{equation}
We shall use the following estimate, proved in Appendix~\ref{app:entropy}.
\begin{lemma}\label{lem:entropy}
For every positive triple,
\begin{equation}\label{eq:entropy-bound}
 0<m\le S/3,\qquad \K\ge\frac{16}{15(S+m)}.
\end{equation}
Equality in $m\le S/3$, or in the lower bound for $\K$, holds if and only
if $x=y=z$.
\end{lemma}

Let $f_t(u)=u(1-u)^t$ and $F_t(E,P)=f_t(x)+f_t(y)+f_t(z)$. For
$\phi(E)=F_t(E,P(E))$, the chain rule and~\eqref{eq:curve} give
\begin{equation}\label{eq:variations}
\begin{aligned}
 \phi'&=\mathcal D_t:=(F_t)_E-m(F_t)_P,\\
 \phi''&=\mathcal C_t:=(F_t)_{EE}-2m(F_t)_{EP}
                  +m^2(F_t)_{PP}+\K(F_t)_P.
\end{aligned}
\end{equation}
At a local maximum containing the selected triple, both directions along
this entropy level curve are feasible, whether or not the entropy constraint
is active. Hence $\mathcal D_t=0$ and $\mathcal C_t\le0$. At fixed $E$,
decreasing $P$ decreases $\HH$, because $\HH_P>0$, and is also feasible.
Local maximality therefore gives $(F_t)_P\ge0$.

If $t>1$ and $(F_t)_P=0$, stationarity gives $(F_t)_E=0$ as well. Since
$(E,P)$ are local coordinates on the fixed-sum plane, this forces
$f_t'(x)=f_t'(y)=f_t'(z)$. Rolle's theorem would then give two distinct
zeros of $f_t''$ between $x$ and $z$, whereas
\[
 f_t''(u)=t(1-u)^{t-2}((t+1)u-2)
\]
has only one zero in $(0,1)$. Thus $R_t:=(F_t)_P>0$. We will prove
\begin{equation}\label{eq:target}
 \mathcal D_t=0,\quad R_t>0,\quad S\le1
 \quad\Longrightarrow\quad \mathcal C_t>0,
\end{equation}
contradicting $\phi''\le0$. The case $t=1$ is immediate:
$F_1=S-S^2+2E$ gives $\mathcal D_1=2$, contradicting stationarity.

\section{The finite coefficient comparison}
For three nonnegative variables $r=(r_1,r_2,r_3)$, let
\[
 G_r(w)=\prod_{i=1}^3(1-r_iw)^{-1}
       =\sum_{j\ge0}h_j(r)w^j,\qquad h_j=0\quad(j<0).
\]
Thus $h_j$ is the complete homogeneous polynomial of degree $j$. We use
\begin{equation}\label{eq:turan}
 h_j(r)^2-h_{j-1}(r)h_{j+1}(r)
 =h_j(r_1r_2,r_1r_3,r_2r_3)\ge0\qquad(j\ge0).
\end{equation}
For distinct positive variables, this identity follows by substituting
$h_j(r)=\sum_i r_i^{j+2}/\prod_{k\ne i}(r_i-r_k)$; continuity gives the
general case. If $r\ne0$, then every $h_j(r)>0$ for $j\ge0$, and
$h_j/h_{j-1}$ is nonincreasing for $j\ge1$, strictly decreasing when all
three variables are positive.

\begin{lemma}\label{lem:coefficient-extremum}
If $r\in[0,1]^3$ and $\sigma=3-\sum_i r_i>0$, then
\begin{equation}\label{eq:b-bound}
 b_j:=h_j(r)-h_{j-1}(r)\le C_\sigma:=\max\{1,\sigma^{-1}\}
 \qquad(j\ge0).
\end{equation}
\end{lemma}
\begin{proof}
The case $j=0$ is immediate. Fix $N\ge1$. If $b_N\le0$, there is nothing
to prove, so suppose $b_N>0$. This excludes $r=0$. By~\eqref{eq:turan},
\[
 \frac{h_j}{h_{j-1}}\ge\frac{h_N}{h_{N-1}}>1
 \qquad(1\le j\le N),
\]
so $b_j>0$ for $0\le j\le N$.

The regime $\sigma\ge2$ is incompatible with $b_N>0$. Indeed, every
monomial in $h_N$ occurs at least once in $(r_1+r_2+r_3)h_{N-1}$, so
\[
 h_N\le(r_1+r_2+r_3)h_{N-1}\le h_{N-1}
 \qquad\text{if }\sigma\ge2.
\]
We may therefore assume $0<\sigma<2$.

For $N\ge2$, differentiation of the generating function gives
\[
 (\partial_{r_i}-\partial_{r_k})b_N
 =(r_i-r_k)[w^{N-2}]
 \frac{(1-w)G_r(w)}{(1-r_iw)(1-r_kw)}.
\]
The coefficient on the right is
\[
 \sum_{\nu=0}^{N-2}h_\nu(r_i,r_k)b_{N-2-\nu},
 \qquad h_\nu(r_i,r_k)=\sum_{a+b=\nu}r_i^a r_k^b,
\]
which is positive whenever $b_N>0$. Thus the derivative is nonnegative
if $r_i\ge r_k$. For $N=1$, $b_1=\sum_i r_i-1$ is constant at fixed sum.

Choose two coordinates in $(0,1)$ with $r_i\ge r_k$, and transfer mass
from the smaller to the larger:
\[
 r_i(s)=r_i+s,\qquad r_k(s)=r_k-s,\qquad
 0\le s\le\min\{1-r_i,r_k\}.
\]
The ordering is preserved since $r_i(s)-r_k(s)=r_i-r_k+2s\ge0$.
Writing $g(s)=b_N(r(s))$, we have $g'(s)\ge0$ whenever $g(s)>0$.
If $g$ had a first zero $s_*$, it would be nondecreasing on $[0,s_*)$,
and continuity would give $g(s_*)\ge g(0)>0$, a contradiction. Hence
positivity persists and $b_N$ does not decrease throughout the transfer.

Keep boundary coordinates fixed and repeat. Each transfer puts at least
one additional coordinate at $0$ or $1$, so at most two transfers leave
at most one coordinate in $(0,1)$. At the fixed sum, the resulting vector
is, up to permutation,
\[
 (1,1,1-\sigma)\quad(0<\sigma\le1),\qquad
 (1,2-\sigma,0)\quad(1\le\sigma<2).
\]
The corresponding coefficients of $(1-w)G_r(w)$ are
\[
 \sum_{k=0}^N(1-\sigma)^k\le\sigma^{-1},\qquad
 (2-\sigma)^N\le1,
\]
respectively. Since the transfers do not decrease the original positive
coefficient, these endpoint bounds prove~\eqref{eq:b-bound}. The formulas
also hold at $\sigma=1$ by continuity.
\end{proof}

\begin{lemma}\label{lem:comparison}
Let $\beta>q_1>q_2>q_3>0$, $\alpha>0$, and $n\ge0$ be an integer.
Write $G=G_q$, $k_j=h_j(q)-\alpha h_{j-1}(q)$, and suppose
$k_n>0$, $k_{n+1}\ge\beta k_n$. Then
\begin{equation}\label{eq:coefficient}
\begin{aligned}
 J&:=[w^{n-1}](1-\alpha w)(1-\beta w)^2G(w)^2\\
  &<\max\left\{\frac1\beta,\frac1{3\beta-\sum_iq_i}\right\}k_n,
\end{aligned}
\end{equation}
where negative-index coefficients are zero.
\end{lemma}
\begin{proof}
For $n=0$, $J=0$, $k_0=1$, and the right-hand side is positive since
$3\beta-\sum_iq_i>0$. Let $n\ge1$, abbreviate $h_j=h_j(q)$, and put
\[
 B_j=h_j-\beta h_{j-1},\qquad
 a_j=k_j-\beta k_{j-1}=B_j-\alpha B_{j-1}.
\]
By~\eqref{eq:turan},
\[
 \frac{h_{n+1}}{h_n}-\frac{k_{n+1}}{k_n}
 =\frac{\alpha(h_n^2-h_{n-1}h_{n+1})}{h_nk_n}>0.
\]
Thus $h_{n+1}/h_n>\beta$, and the decreasing $h$-ratios give
$B_j>0$ for $0\le j\le n+1$.

The partial-fraction representation is
\[
 B_j=\sum_{i=1}^3c_iq_i^j,\qquad
 c_i=\frac{q_i(q_i-\beta)}{\prod_{k\ne i}(q_i-q_k)},
\]
with $c_1<0$, $c_2>0$, $c_3<0$. Define the exponential interpolation
\[
 g(s)=c_1(q_1/q_2)^s+c_2+c_3(q_3/q_2)^s,
 \qquad g(j)=B_j/q_2^j.
\]
We have $g''<0$. Since $g(0)>0$ and $g(n+1)>0$, concavity makes $g$
positive on $[0,n+1]$. Therefore
$(\log g)''=(g''g-(g')^2)/g^2<0$, and the ratios $B_j/B_{j-1}$ strictly
decrease for $1\le j\le n+1$. The hypothesis gives
\[
 B_{n+1}-\alpha B_n=a_{n+1}=k_{n+1}-\beta k_n\ge0.
\]
Consequently $B_j/B_{j-1}>B_{n+1}/B_n\ge\alpha$ for $1\le j\le n$,
so $a_j>0$ for $0\le j\le n$, with $a_0=1$.

Let $C=\max\{1,\beta/(3\beta-\sum_iq_i)\}$. Applying
Lemma~\ref{lem:coefficient-extremum} to $r_i=q_i/\beta$ gives
$B_j\le C\beta^j$ for every $j\ge0$. Hence
\[
 J=\sum_{j=0}^{n-1}a_jB_{n-1-j}
 \le C\sum_{j=0}^{n-1}a_j\beta^{n-1-j}
 =Ck_{n-1}<\frac C\beta k_n.
\]
The equality telescopes using $a_j=k_j-\beta k_{j-1}$ and $k_{-1}=0$;
the last inequality is $a_n>0$.
\end{proof}

\section{Completion of the two-size proof}
\begin{proof}[Proof of Theorem~\ref{thm:main}]
By the three-coordinate reduction, it remains to prove~\eqref{eq:target}
for $t\ge2$. Set $n=t-2$, $\alpha=t/(t+1)$, $\beta=1+m$, and
$q=(1-x,1-y,1-z)$. In the notation of Lemma~\ref{lem:comparison},
\begin{equation}\label{eq:derivatives}
\begin{aligned}
 R_t&=(t+1)k_n,\\
 \mathcal D_t&=(t+1)(k_{n+1}-\beta k_n),\\
 \mathcal C_t&=(t+1)(\K k_n-J).
\end{aligned}
\end{equation}
To verify these identities, denote the second divided difference by
\[
 \psi[x,y,z]=\sum_{u\in\{x,y,z\}}\frac{\psi(u)}{W'(u)}.
\]
Root differentiation gives
$(F_t)_P=f_t'[x,y,z]$ and $(F_t)_E=-(uf_t'(u))[x,y,z]$.
For integer $j\ge0$, the second divided difference of $(1-u)^j$ at
$x,y,z$ is $h_{j-2}(q)$, by partial fractions. Since
\[
 f_t'(u)=(t+1)(1-u)^t-t(1-u)^{t-1},
\]
we obtain $(F_t)_P=(t+1)k_n$ and $(F_t)_E=(t+1)(k_{n+1}-k_n)$,
which give the first two identities in~\eqref{eq:derivatives}.
This is where integrality of $t$ enters: the powers are polynomials,
and $n=t-2$ is a nonnegative integer indexing the coefficient comparison.

For the third identity, observe that
\[
 G_q(w)^{-1}=1-(3-S)w+(3-2S+E)w^2-(1-S+E-P)w^3.
\]
With $m$ frozen at the point of differentiation,
\[
 (\partial_E-m\partial_P)G_q=-w^2(1-\beta w)G_q^2.
\]
Along the entropy curve, $\frac{d}{dE}G_q=(\partial_E-m\partial_P)G_q$
and $\beta'=m'=-\K$. Differentiating
\[
 \frac{\mathcal D_t}{t+1}
 =[w^{n+1}](1-\alpha w)(1-\beta w)G_q
\]
therefore gives
\[
\begin{aligned}
 \frac{\mathcal C_t}{t+1}
 &=[w^{n+1}](1-\alpha w)
       \bigl(\K wG_q-w^2(1-\beta w)^2G_q^2\bigr)\\
 &=\K k_n-J.
\end{aligned}
\]
Here $\mathcal C_t=d\mathcal D_t/dE$ by~\eqref{eq:variations}; the
$\K$ term accounts for the variation of $m$ along the curve.

At a stationary triple, $0<x<y<z<1$ gives
$\beta>q_1>q_2>q_3>0$. Also $R_t>0$ and $\mathcal D_t=0$ give
$k_n>0$ and $k_{n+1}=\beta k_n$. Thus all hypotheses of
Lemma~\ref{lem:comparison} hold. Since $S\le1$, it yields
\[
 \frac J{k_n}<\max\left\{\frac1{1+m},\frac1{S+3m}\right\}
 \le\frac1{S+m}.
\]
Lemma~\ref{lem:entropy} and~\eqref{eq:derivatives} now imply
\begin{equation}\label{eq:positive-curvature}
 \frac{\mathcal C_t}{R_t}=\K-\frac J{k_n}
 >\frac1{15(S+m)}>0.
\end{equation}
This proves~\eqref{eq:target}, excluding every local maximum with three
distinct positive coordinate values. The same calculation also gives the
quantitative bound $\mathcal C_t>R_t/(20S)$, since $m\le S/3$.

These perturbations change only three coordinates and tend to zero in
$\ell^1$, so the argument applies equally to countably infinite alphabets.
Each fixed positive size can occur only finitely often in a probability
vector. Hence at most two positive sizes imply finite support.

For completeness, we give the compactness argument
of~\cite[Lemmas~1--2 and Theorem~2]{BKZ}. On a countably infinite alphabet,
identify $I$ with $\mathbb N$. Both $H$ and $\Phi_t$ are symmetric and are
unaffected by insertion or deletion of zero coordinates. It therefore
suffices to take a maximizing sequence whose coordinates are nonincreasing.
For such a probability vector, $p_i\le1/i$, and for $N\ge1$,
\[
 h\ge H(p)\ge\sum_{i>N}p_i\log\frac1{p_i}
 \ge\log(N+1)\sum_{i>N}p_i.
\]
Choose a coordinatewise convergent subsequence $p^{(k)}\to p$ by a diagonal
argument. The uniform tail bound prevents loss of mass: passing to the limit
in finite partial sums and then letting $N\to\infty$ gives
$\sum_i p_i=1$. The same tail bound, together with
convergence on each finite initial segment, gives
$\|p^{(k)}-p\|_1\to0$. Since $-u\log u\ge0$ on $[0,1]$, Fatou's lemma
gives
\[
 H(p)\le\liminf_{k\to\infty}H(p^{(k)})\le h.
\]
Finally, $f_t\in C^1[0,1]$ and $f_t(0)=0$, so
\[
 |\Phi_t(p)-\Phi_t(q)|\le L_t\|p-q\|_1,
 \qquad L_t=\max_{0\le u\le1}|f_t'(u)|<\infty.
\]
Thus the maximum is attained. On a finite alphabet, attainment follows
directly from compactness of the feasible set and continuity of the objective.

For entropy saturation, use the atom-splitting argument
of~\cite[Proposition~1]{BKZ}. Let $I$ be countably infinite and let $p$ be
a local maximizer with $H(p)<h$. Its support is finite, so an unused
coordinate is available. Replace a positive atom $a$ by $a-\varepsilon$
and put $\varepsilon$ in that unused coordinate, where $0<\varepsilon<a$.
The new vector $p^{\varepsilon}$ satisfies
\[
 \|p^{\varepsilon}-p\|_1=2\varepsilon,\qquad
 H(p^{\varepsilon})-H(p)
 =a\log a-(a-\varepsilon)\log(a-\varepsilon)
                -\varepsilon\log\varepsilon\longrightarrow0.
\]
The entropy increment is positive, and
\[
 f_t(a-\varepsilon)+f_t(\varepsilon)
 >(a-\varepsilon)(1-a)^t+\varepsilon(1-a)^t=f_t(a).
\]
For sufficiently small $\varepsilon$, this is a feasible improvement
arbitrarily close to $p$, contradicting local maximality. Hence every
local maximizer satisfies $H(p)=h$.
\end{proof}

\section{One exceptional atom and a finite classification}\label{sec:classification}

\subsection{The multiplicity restriction}
The following elementary criterion isolates the additional second-variation
argument. It uses only perturbations of four coordinates.

\begin{lemma}[Repeated-size criterion]\label{lem:repeated}
Let $f\in C[0,1]\cap C^2(0,1)$, and let $J\subset(0,1)$ be an interval
containing every $u$ for which $f''(u)\le0$. Suppose that $g(u)=uf''(u)$
satisfies
\begin{equation}\label{eq:strict-quasiconvex}
 g(v)<\max\{g(x),g(y)\}\qquad(x<v<y,\quad x,y\in J).
\end{equation}
At a finite-dimensional local maximum of $\sum_i f(p_i)$ under normalization
and an upper Shannon-entropy constraint, at most one distinct positive
coordinate value can occur more than once. The same conclusion holds on a
countable alphabet whenever the objective is well-defined and the point is
locally maximal under finite-coordinate perturbations.
\end{lemma}
\begin{proof}
If a value $u$ occurs twice, replacing these copies by $u+\varepsilon$ and
$u-\varepsilon$ preserves mass and strictly decreases entropy. Local
maximality and Taylor expansion therefore imply $f''(u)\le0$.

Suppose two distinct values $x<y$ both occur at least twice, and vary just
$(x,x,y,y)$. On these four coordinates, the gradients of mass and entropy
are independent. The fixed-mass, fixed-entropy set through the point is
therefore a smooth manifold, whether or not the entropy inequality is
active. Equality-constrained first-order conditions give constants
$\lambda,\mu$ such that
\[
 f'(u)=\lambda+\mu(-\log u-1),\qquad u\in\{x,y\}.
\]
In particular,
\begin{equation}\label{eq:log-average}
 -\mu=\frac{f'(y)-f'(x)}{\log(y/x)}
 =\frac{\displaystyle\int_x^y g(u)\,\frac{\dd u}{u}}
        {\log(y/x)}.
\end{equation}
A pair-difference direction within either repeated block is tangent to both
constraints. The second-order condition on this manifold is
\[
 f''(u)+\frac\mu u\le0,\qquad u\in\{x,y\}.
\]
Thus $g(x),g(y)\le-\mu$. But $x,y\in J$, and
\eqref{eq:strict-quasiconvex} makes the weighted average
in~\eqref{eq:log-average} strictly smaller than
$\max\{g(x),g(y)\}$, a contradiction. All perturbations used here have
finite support, which also proves the countable-alphabet assertion.
\end{proof}

\begin{theorem}[One exceptional atom]\label{thm:exceptional}
Under the hypotheses of Theorem~\ref{thm:main}, every nonuniform local
maximizer has, up to permutation and zero coordinates, the form
\begin{equation}\label{eq:exceptional-form}
 p_m(z)=\left(z,\underbrace{q,\ldots,q}_{m\text{ copies}}\right),
 \qquad q=\frac{1-z}{m},\quad m\in\mathbb N,\quad 0<z<1.
\end{equation}
Thus either $z<q$ (one light atom) or $z>q$ (one heavy atom).
\end{theorem}
\begin{proof}
For $t=1$, $u f_t''(u)=-2u$ is strictly decreasing. For $t>1$, the set
where $f_t''\le0$ is $J=(0,2/(t+1)]$, and
\[
 \frac{\dd}{\dd u}\bigl(uf_t''(u)\bigr)
 =-t(1-u)^{t-3}\bigl(t(t+1)u^2-4tu+2\bigr).
\]
The quadratic roots are
\[
 r_\pm=\frac{2t\pm\sqrt{2t(t-1)}}{t(t+1)}.
\]
They satisfy $0<r_-<2/(t+1)<r_+$. Hence $uf_t''(u)$ strictly decreases
and then strictly increases on $J$, so it satisfies
\eqref{eq:strict-quasiconvex}. Lemma~\ref{lem:repeated} shows that two
distinct repeated sizes are impossible. Theorem~\ref{thm:main} supplies
finite support and at most two sizes, so at least one of their
multiplicities is one.
\end{proof}

A uniform law on $m+1$ points is the limiting case $z=q=1/(m+1)$.
A nonuniform binary law has both a one-light and a one-heavy description,
which differ only by permutation. For support size at least three the two
nonuniform descriptions are distinct.

\subsection{The entropy-matched candidates}
For the rest of this section the alphabet is countably infinite. Define
\[
 B_t(h)=\max_{p\in\mathcal P_h(\mathbb N)}\Phi_t(p).
\]
If $h=0$, only point masses are feasible, and $B_t(0)=0$. Fix $h>0$ and put
$k=\lfloor e^h\rfloor$. For $m\ge1$, set
\begin{equation}\label{eq:branch-functions}
 \begin{aligned}
 E_m(z)&=-z\log z-(1-z)\log\frac{1-z}{m},\\
 V_{m,s}(z)&=z(1-z)^s+(1-z)\left(1-\frac{1-z}{m}\right)^s.
 \end{aligned}
\end{equation}
The functions extend continuously to the entropy endpoints. Since
\begin{equation}\label{eq:entropy-branch-derivative}
 E_m'(z)=\log\frac{1-z}{mz},
\end{equation}
$E_m$ increases from $\log m$ to $\log(m+1)$ on $[0,1/(m+1)]$ and
decreases from $\log(m+1)$ to $0$ on $[1/(m+1),1]$.

There is a unique $a\in[0,1/(k+1))$ with $E_k(a)=h$. Define
\begin{equation}\label{eq:light-candidate}
 p_{\rm L}=p_k(a),\qquad y=\frac{1-a}{k},\qquad
 \ell(s)=V_{k,s}(a).
\end{equation}
If $h=\log k$, then $a=0$, and deleting its zero coordinate gives the
uniform law on $k$ points. Otherwise $p_{\rm L}$ is one-light and has
$k+1$ positive coordinates.

For every integer $m\ge k$, there is a unique
$z_m\in(1/(m+1),1)$ with $E_m(z_m)=h$. Define
\begin{equation}\label{eq:heavy-candidates}
 q_m=\frac{1-z_m}{m},\qquad
 p_m^{\rm H}=p_m(z_m),\qquad b_m(s)=V_{m,s}(z_m).
\end{equation}
The dependence on the fixed entropy $h$ is suppressed in this notation.
Theorem~\ref{thm:exceptional} and entropy saturation show that every global
maximizer is among $p_{\rm L}$ and the $p_m^{\rm H}$, $m\ge k$.
When $0<h<\log2$, $p_1^{\rm H}$ is the binary duplicate of $p_{\rm L}$.
The real parameter $s$ in $\ell(s)$ and $b_m(s)$ will also be useful for
comparing these candidate values; the structural theorems continue to
require an integer sample size.

\subsection{A uniform cutoff for the heavy family}
\begin{lemma}[Heavy-tail cutoff]\label{lem:cutoff}
Fix $h>0$ and a real exponent $s\ge1$. Interpolate the heavy formulas by
allowing $m>e^h-1$ to be real and taking the unique root
$E_m(z_m)=h$ with $z_m>1/(m+1)$. Then
\begin{equation}\label{eq:cutoff-decrease}
 \frac{\dd}{\dd m}b_m(s)<0\qquad\text{whenever }m\ge sh+1.
\end{equation}
\end{lemma}
\begin{proof}
This is an interpolation of the formulas, not a probability vector with
nonintegral multiplicities. Write $z=z_m$, $q=(1-z)/m$, and
$\sigma=\log(z/q)>0$. The entropy equation is
\begin{equation}\label{eq:heavy-entropy-sigma}
 h=-\log z+(1-z)\sigma.
\end{equation}
Implicit differentiation gives
\begin{equation}\label{eq:heavy-m-derivatives}
 \frac{\dd z}{\dd m}=\frac q\sigma,
 \qquad
 \frac{\dd q}{\dd m}=-\frac{q(1+\sigma)}{m\sigma}.
\end{equation}
With $f_s(u)=u(1-u)^s$, it follows that
\begin{equation}\label{eq:heavy-objective-derivative}
 \frac{\dd b_m(s)}{\dd m}=\frac q\sigma\Psi_s(z,q),\qquad
 \Psi_s(z,q)=f_s'(z)-f_s'(q)+sq\sigma(1-q)^{s-1}.
\end{equation}

If $z\le1/s$, hold $z$ fixed and put
\[
 A_z(u)=f_s'(u)-su\log(z/u)(1-u)^{s-1}.
\]
Direct differentiation yields
\[
 A_z'(u)=s\bigl(1+\log(z/u)\bigr)(1-u)^{s-2}(su-1).
\]
This is strictly negative for $q<u<z$. Hence
$\Psi_s(z,q)=A_z(z)-A_z(q)<0$.

If $z>1/s$, then $f_s'(z)<0$, and
\[
 \Psi_s(z,q)=f_s'(z)
 -(1-q)^{s-1}\bigl(1-q-sq(1+\sigma)\bigr).
\]
By~\eqref{eq:heavy-entropy-sigma} and $\log z<z-1$,
\[
 (1-z)(1+\sigma)=h+\log z+1-z<h.
\]
Consequently $sq(1+\sigma)<sh/m$ and $q<1/m$. If $m\ge sh+1$, then
\[
 1-q-sq(1+\sigma)>1-\frac{sh+1}{m}\ge0.
\]
Thus $\Psi_s(z,q)<0$ in this case as well.
\end{proof}

\begin{theorem}[Exact finite classification]\label{thm:finite-classification}
For integer $t\ge1$ and $h>0$, put
\begin{equation}\label{eq:M-cutoff}
 M=M(t,h)=\max\{k,\lceil th+1\rceil\},\qquad k=\lfloor e^h\rfloor.
\end{equation}
Then
\begin{equation}\label{eq:finite-maximum}
 B_t(h)=\max\{\ell(t),b_k(t),b_{k+1}(t),\ldots,b_M(t)\}.
\end{equation}
Up to permutation and zero coordinates, the global maximizers are exactly
the candidates in this list attaining the displayed maximum. In particular,
every global maximizer has at most $M+1$ positive coordinates.
\end{theorem}
\begin{proof}
The candidate list before truncation is exhaustive by
Theorem~\ref{thm:exceptional} and entropy saturation. Every candidate is
feasible. Lemma~\ref{lem:cutoff} gives $b_j(t)<b_M(t)$ for every integer
$j>M$, which proves the formula and its equality statement.
\end{proof}

Thus the atom sizes are specified by unique monotone entropy roots, and the
remaining comparison is finite. All ties are retained. Appendix~\ref{app:crossings}
gives crossing intervals specifying when each candidate wins;
Appendix~\ref{app:finite-alphabet} treats a fixed finite alphabet, where the
entropy constraint can be slack.

\subsection{Sample sizes one and two}
\begin{corollary}\label{cor:small-samples}
For $t=1$ and $t=2$, the unique global maximizer on a countably infinite
alphabet, up to permutation and zero coordinates, is $p_{\rm L}$.
Equivalently, $B_t(h)=\ell(t)$ for $h>0$.
\end{corollary}
The $t=1$ case recovers the maximizing shape in
\cite[Proposition~3]{BKZ}.
\begin{proof}
It suffices to exclude every distinct heavy candidate with $m\ge2$.
Write $z=rq$, $q=1/(m+r)$, and $r>1$. The fixed-mass, fixed-entropy
second-order condition in a pair-difference direction among the repeated
$q$-atoms is
\begin{equation}\label{eq:small-sample-hessian}
 qf_t''(q)+\mu\le0,\qquad
 \mu=\frac{f_t'(q)-f_t'(z)}{\log r},
\end{equation}
by the same equality-manifold argument as in Lemma~\ref{lem:repeated}.
For $t=1$, the left side is
\[
 2q\left(\frac{r-1}{\log r}-1\right)>0.
\]
For $t=2$, it is
\[
 \frac{N_m(r)}{(m+r)^2\log r},\qquad
 N_m(r)=(r-1)(4m+r-3)-(4m+4r-6)\log r.
\]
Now
\[
 \begin{aligned}
 N_m(r)&=N_2(r)+4(m-2)(r-1-\log r),\\
 N_2(r)&=(r-1)(r+5)-(4r+2)\log r,\\
 N_2'(r)&=2\left(r-\frac1r-2\log r\right)>0\qquad(r>1).
 \end{aligned}
\]
The last inequality follows by differentiating its parenthesis, whose
derivative is $(r-1)^2/r^2$ and whose value at $1$ is zero. Since
$N_2(1)=0$, we have $N_m(r)>0$, again contradicting
\eqref{eq:small-sample-hessian}. The only possible heavy candidate with
$m=1$ is the binary duplicate of $p_{\rm L}$. Attainment now proves the
claim, including uniqueness.
\end{proof}

\section{Sharp asymptotics and occupancy consequences}\label{sec:asymptotics}
Throughout this section the entropy $h>0$ is fixed and the alphabet is
countably infinite. Berend, Kontorovich, and Zagdanski proved
$B_t(h)\sim h/\log t$ by matching first-order upper and lower
bounds~\cite[Theorems~3--4]{BKZ}. The classification yields the following
refinement and the scales of all optimizing atoms.

\begin{theorem}[Fixed-entropy asymptotics]\label{thm:asymptotics}
As integer $t\to\infty$,
\begin{equation}\label{eq:sharp-asymptotic}
 \frac{h}{B_t(h)}
 =\log t+\log\log t+2
  +O_h\!\left(\frac{\log\log t}{\log t}\right).
\end{equation}
Every global maximizer is eventually one-heavy. Writing any such maximizer
as
\[
 p_t=\left(1-L_t,
       \underbrace{q_t,\ldots,q_t}_{m_t\text{ copies}}\right),
 \qquad m_tq_t=L_t,
\]
we have, uniformly over the choice of a maximizer,
\begin{equation}\label{eq:optimizer-scales}
 L_t\sim\frac h{\log t},\qquad
 q_t\sim\frac1{t\log t},\qquad
 m_t\sim ht.
\end{equation}
\end{theorem}
\begin{proof}
Put $T=\log t$. For a heavy vector with total light mass $L=mq$, the exact
entropy and objective are
\begin{equation}\label{eq:L-entropy-objective}
 \begin{aligned}
 h&=-(1-L)\log(1-L)-L\log q,\\
 \Phi_t(p)&=(1-L)L^t+L(1-q)^t.
 \end{aligned}
\end{equation}
It is convenient to set
\begin{equation}\label{eq:A-L}
 A(L)=-\frac{1-L}{L}\log(1-L)
     =1-\frac L2+O(L^2)\qquad(L\downarrow0).
\end{equation}
All constants below may depend on $h$, but not on the maximizing vector.

\paragraph{A lower bound with integer multiplicity.}
Choose $m=\lceil ht\rceil$ and solve the heavy entropy equation for $L$.
For sufficiently large $t$, $h<\log(m+1)$, so this root exists. Its entropy
can also be written as $h=H_2(L)+L\log m$, where
$H_2(L)=-L\log L-(1-L)\log(1-L)\ge0$. Thus
$L\le h/\log m=O_h(T^{-1})$. With $d=h/L$, the entropy equation becomes
\[
 d=\log(m/h)+\log d+A(L)
  =T+\log d+1+O_h(T^{-1}).
\]
Since $d\ge\log m$ and $d-\log d$ is increasing for $d>1$, this gives
\begin{equation}\label{eq:d-lower-candidate}
 d=T+\log T+1+O_h\!\left(\frac{\log T}{T}\right).
\end{equation}
For example, substituting $T+\log T+1$ in $d-\log d$ gives an error
$O(\log T/T)$, and the derivative is bounded away from zero in this range.
Moreover,
\[
 tq=\frac{th}{md}=\frac1d+O_h\!\left(\frac1{tT}\right),
 \qquad
 -t\log(1-q)=\frac1d+O_h\!\left(\frac1{tT}\right).
\]
Since $B_t(h)\ge L(1-q)^t$, it follows that
\begin{equation}\label{eq:reciprocal-upper}
 \begin{aligned}
 \frac h{B_t(h)}
 &\le d\exp\!\left(\frac1d+O_h\!\left(\frac1{tT}\right)\right)\\
 &=T+\log T+2+O_h\!\left(\frac{\log T}{T}\right).
 \end{aligned}
\end{equation}
In particular $B_t(h)\ge c_h/T$ for some $c_h>0$ and all sufficiently
large $t$.

\paragraph{A uniform upper bound.}
At fixed $h$, the light or uniform candidate has a fixed finite positive
support, so its objective decays exponentially in $t$. The preceding
lower bound and Theorem~\ref{thm:finite-classification} therefore imply
that every maximizer is eventually heavy.

Consider any such maximizer, with largest atom $z=1-L$. Since
$H(p)\ge-\log z$, we have $z\ge e^{-h}$ and
$L\le\rho_h:=1-e^{-h}<1$. The heavy atom contributes at most $\rho_h^t$.
Put $u=tq$ and $d=h/L$. From
\[
 L(1-q)^t=B_t(h)-(1-L)L^t\ge\frac{c_h}{2T}
\]
for large $t$, and $(1-q)^t\le e^{-u}$, we get
$u\le\log(2T/c_h)=O_h(\log T)$. The entropy identity gives
$h\ge L(T-\log u)$, so $L=O_h(T^{-1})$. Therefore
\begin{equation}\label{eq:optimizer-d}
 d=T-\log u+1+O_h(T^{-1}).
\end{equation}
Let $R=Le^{-u}$. We have $R\ge c_h/(2T)$ and
$B_t(h)\le R+\rho_h^t$. Consequently
\begin{equation}\label{eq:reciprocal-uniform}
 \frac h{B_t(h)}
 \ge de^u-O_h(T^2\rho_h^t)
 =de^u-o_h(T^{-1}).
\end{equation}
This explicitly controls the error from discarding the heavy atom.

The bounds~\eqref{eq:reciprocal-upper} and
\eqref{eq:reciprocal-uniform} force $u<1$ for all sufficiently large $t$.
Indeed, if $u\ge1$, then~\eqref{eq:optimizer-d} and
$u\le\log(2T/c_h)$ imply
\[
 de^u\ge e\bigl(T-O_h(\log\log T)\bigr),
\]
contradicting~\eqref{eq:reciprocal-upper}. For $u<1$, equation
\eqref{eq:optimizer-d} gives $d\ge T$ for large $t$. Using
$e^u\ge1+u$ and putting $v=Tu$, we obtain
\begin{equation}\label{eq:scalar-gap}
 \begin{aligned}
 de^u&\ge d+Tu\\
 &\ge T+\log T+1+v-\log v-O_h(T^{-1})\\
 &\ge T+\log T+2-O_h(T^{-1}),
 \end{aligned}
\end{equation}
because $v-\log v\ge1$. Combining this with
\eqref{eq:reciprocal-upper} and~\eqref{eq:reciprocal-uniform} proves
\eqref{eq:sharp-asymptotic}.

\paragraph{The optimizing scales.}
The same inequalities give
\[
 0\le v-\log v-1
 =O_h\!\left(\frac{\log T}{T}\right).
\]
The function on the left has its unique zero at $v=1$ and tends to infinity
at both ends of $(0,\infty)$. Hence $v\to1$, uniformly over all maximizers.
Thus $tq=u\sim1/T$. Equation~\eqref{eq:optimizer-d} now gives $d\sim T$,
so $L=h/d\sim h/T$ and $m=L/q\sim ht$, as required.
\end{proof}

\begin{corollary}[Expectation-level sample complexity]\label{cor:sample-complexity}
For fixed $h>0$, define
\[
 N_h(\varepsilon)=\min\{t\in\mathbb N:B_t(h)\le\varepsilon\}.
\]
Then
\begin{equation}\label{eq:sample-complexity}
 N_h(\varepsilon)\sim
 \frac{\varepsilon}{e^2h}\exp\!\left(\frac h\varepsilon\right)
 \qquad(\varepsilon\downarrow0).
\end{equation}
\end{corollary}
\begin{proof}
The function $B_t(h)$ is strictly decreasing in integer $t$: at an optimizer
for $t+1$, positive entropy implies that every positive atom is below one,
so $\Phi_t(p)>\Phi_{t+1}(p)$. Also $B_t(h)\to0$ by
Theorem~\ref{thm:asymptotics}. Put $A=h/\varepsilon$ and
$S=A-\log A-2$. For any fixed $\eta>0$, apply
\eqref{eq:sharp-asymptotic} at the integers nearest $e^{S-\eta}$ and
$e^{S+\eta}$. Their reciprocal objective values are
$A-\eta+o(1)$ and $A+\eta+o(1)$, respectively. They therefore bracket
$N_h(\varepsilon)$. It follows that
$\log N_h(\varepsilon)=S+o(1)$, proving the assertion.
\end{proof}

\begin{corollary}[Singletons, discovery, and coverage]\label{cor:occupancy}
Let $K_n$ be the number of distinct symbols in $n$ independent samples,
let $K_{n,1}$ be the number occurring exactly once, and let $C_n$ be the
probability mass of the observed symbols. Then
\begin{equation}\label{eq:occupancy-extrema}
 \begin{aligned}
 \max_{H(p)\le h}\EE_p K_{n,1}&=nB_{n-1}(h) &&(n\ge2),\\
 \max_{H(p)\le h}\EE_p(K_{t+1}-K_t)&=B_t(h) &&(t\ge1),\\
 \min_{H(p)\le h}\EE_p C_t&=1-B_t(h) &&(t\ge1).
 \end{aligned}
\end{equation}
The optimizers are exactly those in Theorem~\ref{thm:finite-classification},
with the corresponding sample-size index.
\end{corollary}
\begin{proof}
Linearity of expectation gives
\[
 \EE_p K_{n,1}=n\sum_i p_i(1-p_i)^{n-1},\qquad
 \EE_p(K_{t+1}-K_t)=\sum_i p_i(1-p_i)^t,
\]
and $\EE_p C_t=1-\Phi_t(p)$. These identities prove every assertion.
\end{proof}
In particular the singleton estimator $K_{n,1}/n$ has maximal expectation
$B_{n-1}(h)$. The optimization in~\eqref{eq:occupancy-extrema} concerns
expectations, rather than estimation error or tail probabilities.

\appendix
\section{The entropy-curvature estimate}\label{app:entropy}
\begin{proof}[Proof of Lemma~\ref{lem:entropy}]
For a uniform triple $x=y=z=a$, direct integration gives
\[
 I_0=\frac1{2a^2},\qquad I_1=\frac1{2a},\qquad
 m=a,\qquad \K=\frac4{15a}=\frac{16}{15(S+m)}.
\]
This proves the lemma, with equality, in the uniform case. We first establish
a kernel estimate and its mixture form, and then apply them to a nonuniform
triple.

For $a>0$, set $\kappa_a(s)=2a^2/(s+a)^3$, a probability density on
$[0,\infty)$ of mean $a$, and
$J_*(a)=\int_0^\infty(s-1)^2\kappa_1(s)\kappa_a(s)\dd s$.
The kernel estimate is
\begin{equation}\label{eq:kernel}
 J_*(a)\ge\frac{2(9-a)}{15(a+1)}
 =\frac{13}{15}-\frac a3+\frac{(a-1)^2}{3(a+1)}.
\end{equation}
Here is an explicit positive remainder, with $b_v=1+v(a-1)$:
\[
 J_*(a)-\frac{2(9-a)}{15(a+1)}
 =\frac{2(a-1)^2}{5(a+1)}\int_0^1
 \frac{v(1-v)^5(2b_v^3+12b_v^2-10b_v+13)}{b_v^4}\dd v.
\]
The cubic is positive since it equals
$2b_v^3+12(b_v-5/12)^2+131/12$ and $b_v>0$.
For verification, direct integration gives, for $a\ne1$,
\[
 J_*(a)=\frac{4a^2(a^2+10a+13)\log a
       -2(a-1)(7a^3+33a^2+9a-1)}{(a-1)^5}.
\]
For $M(a)=\frac{15}{2}(a-1)^5(a+1)(J_*(a)-2(9-a)/(15(a+1)))$, with
its smooth extension at $a=1$, one has
\[
 M^{(j)}(1)=0\quad(0\le j\le5),\qquad
 M^{(6)}(a)=\frac{360(a-1)(2a^3+12a^2-10a+13)}{a^4}.
\]
Taylor's integral remainder gives the displayed identity on both sides
of $1$; directly $J_*(1)=8/15$.

For a probability law $\lambda$ supported in a compact subset of $(0,\infty)$,
put $m_\lambda=\int v\dd\lambda(v)$ and
$G(s)=\int\kappa_v(s)\dd\lambda(v)$. Scaling~\eqref{eq:kernel} gives
\[
\begin{aligned}
 \int_0^\infty(s-m_\lambda)^2\kappa_{m_\lambda}(s)\kappa_v(s)\dd s
 &=m_\lambda J_*(v/m_\lambda)\\
 &\ge\frac{13m_\lambda}{15}-\frac v3
       +\frac{(v-m_\lambda)^2}{3(v+m_\lambda)}.
\end{aligned}
\]
Averaging and using $\int v\dd\lambda(v)=m_\lambda$ yields
\[
 \int_0^\infty(s-m_\lambda)^2\kappa_{m_\lambda}(s)G(s)\dd s
 \ge\frac{8m_\lambda}{15}
    +\frac13\int\frac{(v-m_\lambda)^2}{v+m_\lambda}\dd\lambda(v).
\]
Expand the nonnegative integral
$\int_0^\infty(s-m_\lambda)^2(G-\kappa_{m_\lambda})^2\dd s$ and use
$\int_0^\infty(s-m_\lambda)^2\kappa_{m_\lambda}(s)^2\dd s=8m_\lambda/15$.
It follows that
\begin{equation}\label{eq:mixture}
 \int_0^\infty(s-m_\lambda)^2G(s)^2\dd s
 \ge\frac{8m_\lambda}{15}
      +\frac23\int\frac{(v-m_\lambda)^2}{v+m_\lambda}\dd\lambda(v).
\end{equation}

Now let $(x,y,z)$ be nonuniform. Let $(U_1,U_2)$ have density $2$ on
$\{(u_1,u_2):u_1,u_2\ge0,\ u_1+u_2\le1\}$, set $U_3=1-U_1-U_2$, and
put $V=xU_1+yU_2+zU_3$. Thus $(U_1,U_2,U_3)$ is uniform on the probability
simplex and $\EE V=S/3$. Two elementary integrations over the simplex give
\[
 D(s)^{-1}=\EE(s+V)^{-3},\qquad
 I_0=\tfrac12\EE V^{-2},\quad I_1=\tfrac12\EE V^{-1}.
\]
Tilt its law by $\dd\lambda(v)=v^{-2}\dd\mathbb P_V(v)/(2I_0)$. Then
\[
 G(s)=\frac1{I_0D(s)},\quad
 \EE_\lambda V=m,\quad b_2:=\EE_\lambda V^2=\frac1{2I_0},\quad
 \EE_\lambda V^3=(S/3)b_2.
\]
Consequently
\[
 \EE_\lambda[(V-m)^2(V+m)]=(S/3-m)b_2.
\]
Since $V$ is positive and nonconstant, the left-hand side is strictly
positive, so $0<m<S/3$. In particular, the following denominators are
positive. Cauchy--Schwarz gives
\[
 \EE_\lambda\frac{(V-m)^2}{V+m}
 \ge\frac{\bigl(\EE_\lambda(V-m)^2\bigr)^2}
            {\EE_\lambda[(V-m)^2(V+m)]}
 =\frac{(b_2-m^2)^2}{(S/3-m)b_2}.
\]
Substitute this bound into~\eqref{eq:mixture}, where $m_\lambda=m$, and
multiply by $I_0$ to obtain
\[
 \K\ge\frac8{15}I_1+\frac{(1-2mI_1)^2}{S-3m}.
\]
Also $D(s)\le(s+S/3)^3$ by the arithmetic--geometric mean inequality,
so $j:=SI_1\ge3/2$. Put $r=m/S$, $d=1-3r>0$, and $v=j-3/2\ge0$.
The identities
\[
 \frac8{15}j+\frac{(1-2rj)^2}{d}-\frac{36-48r}{25}
 =\frac{(2rv-3d/5)^2}{d}+\frac8{15}dv\ge0,
\]
\[
 (1+r)\frac{36-48r}{25}-\frac{16}{15}
 =\frac{4(1-3r)(12r+7)}{75}>0
\]
give $\K>16/(15(S+m))$ for a nonuniform triple. Together with the uniform
calculation, this proves the lemma and its equality assertions.
\end{proof}

\section{Crossing structure of the candidate family}\label{app:crossings}
Fix $h>0$ and retain the notation of Section~\ref{sec:classification}.
For real $s\ge0$, $\ell(s)$ and $b_m(s)$ are exponential sums associated
with fixed finite-support distributions. The following comparisons describe
their crossings exactly. Their use in classifying the original optimization
is restricted to integer sample sizes.

\subsection{Two elementary identities}
\begin{lemma}\label{lem:entropy-series}
For every discrete probability distribution,
\begin{equation}\label{eq:entropy-series}
 H(p)=\sum_{j=1}^\infty\frac{\Phi_j(p)}j,
\end{equation}
with equality in $[0,\infty]$. In particular, if two distributions have
the same finite entropy, the weighted series of the differences of their
$\Phi_j$ values converges absolutely and sums to zero.
\end{lemma}
\begin{proof}
For $0<u\le1$,
$-\log u=\sum_{j\ge1}(1-u)^j/j$. Multiplication by $p_i$ and Tonelli's
theorem prove~\eqref{eq:entropy-series}. Absolute convergence of the
difference follows by bounding it termwise by the sum of the two
nonnegative series.
\end{proof}

\begin{lemma}[Exponential-sum zero counts]\label{lem:zero-count}
A nonzero sum of $r$ exponentials with distinct positive bases has at most
$r-1$ real zeros, counted with multiplicity. If
\[
 F(s)=Aa^s-Bb^s-Cc^s+Dd^s,
 \qquad 0<a<b<c<d,\quad A,B,C,D>0,
\]
then $F$ has at most two real zeros, counted with multiplicity.
\end{lemma}
\begin{proof}
For the first assertion, divide by one exponential and differentiate.
The derivative has $r-1$ exponential terms. Induction and Rolle's theorem,
including repeated zeros, prove the claim.

For the second assertion, put $U(s)=F(s)/b^s$. Then
\[
 \frac{U'(s)}{(d/b)^s}
 =A\log(a/b)(a/d)^s-C\log(c/b)(c/d)^s+D\log(d/b).
\]
The right-hand side has strictly positive derivative, so it has at most
one zero, necessarily simple. Hence $U'$ has at most one zero, counted
with multiplicity. Rolle's theorem gives at most two for $U$ and $F$.
\end{proof}

\subsection{Heavy--heavy and light--heavy crossings}
\begin{theorem}[Heavy--heavy single crossing]\label{thm:heavy-crossing}
For each pair of integers $k\le i<j$, there is a unique positive crossing
$\tau_{ij}>1$. It is simple, and
\begin{equation}\label{eq:heavy-crossing-sign}
 \begin{cases}
 b_j(s)<b_i(s),&0<s<\tau_{ij},\\
 b_j(s)>b_i(s),&s>\tau_{ij}.
 \end{cases}
\end{equation}
\end{theorem}
\begin{proof}
Equations~\eqref{eq:heavy-m-derivatives} give $z_i<z_j$ and $q_j<q_i$.
Since $q_i<z_i$, the bases in
\[
 b_j(s)-b_i(s)
 =z_j(1-z_j)^s-z_i(1-z_i)^s
  -(1-z_i)(1-q_i)^s+(1-z_j)(1-q_j)^s
\]
are strictly increasing in the displayed order, and their coefficient
signs are $+,-,-,+$. Lemma~\ref{lem:zero-count} gives at most two real
zeros, counting multiplicity. One is $s=0$, since both distributions
have mass one. The difference is positive for all sufficiently large $s$,
because $q_j<q_i$. By Lemma~\ref{lem:entropy-series}, it must be negative
at some positive integer. Hence there is exactly one positive crossing,
and the zeros at zero and at that crossing are both simple. The sign
orientation follows, and the negative integer implies $\tau_{ij}>1$.
\end{proof}

\begin{theorem}[Light--heavy crossings]\label{thm:light-crossing}
Fix $m\ge k$ such that $p_m^{\rm H}$ is not the binary duplicate of
$p_{\rm L}$. There are endpoints
$2<\alpha_m<\beta_m\le\infty$ such that
\begin{equation}\label{eq:light-crossing-sign}
 b_m(s)>\ell(s)\quad\Longleftrightarrow\quad
 \alpha_m<s<\beta_m\qquad(s>0).
\end{equation}
More precisely, if $p_{\rm L}$ is uniform or $q_m\le a$, then
$\beta_m=\infty$ and $\alpha_m$ is the unique positive zero of
$b_m-\ell$. If $0<a<q_m$, then $\alpha_m$ and $\beta_m$ are its
exactly two positive zeros. Every finite positive crossing is simple;
outside the closed interval the inequality in~\eqref{eq:light-crossing-sign}
is strictly reversed.
\end{theorem}
\begin{proof}
First note the ordering $q_m<y<z_m$. The first inequality follows from
$q_m<1/(m+1)\le1/(k+1)<y$. To prove $z_m>y$, keep one $y$-atom of
$p_{\rm L}$ and average its remaining masses over $m$ coordinates, padding
with zeros if necessary. This strictly increases entropy unless $k=m=1$,
the excluded binary duplicate. Thus $E_m(y)>h$. Since
$y>1/(m+1)$ and $E_m$ is strictly decreasing on the heavy branch,
$z_m>y$ follows. This averaging argument also applies to the uniform
case $a=0$.

Corollary~\ref{cor:small-samples} gives
\begin{equation}\label{eq:light-strict-small}
 b_m(1)<\ell(1),\qquad b_m(2)<\ell(2).
\end{equation}
If $a=0$, the difference is a three-term exponential sum with increasing
bases and signs $+,-,+$. It is positive for large $s$. By
Lemma~\ref{lem:zero-count}, the zero at $s=0$ and a single simple
crossing $\alpha_m>2$ exhaust its zeros.

Suppose $a>0$. We have
\[
 b_m(s)-\ell(s)
 =z_m(1-z_m)^s-(1-a)(1-y)^s
  +(1-z_m)(1-q_m)^s-a(1-a)^s.
\]
If $q_m<a$, the increasing-base signs are $+,-,-,+$, and the difference
is positive for large $s$. The same two-zero bound gives a unique simple
positive crossing $\alpha_m>2$. If $q_m=a$, the last two terms with
base $1-a$ combine, with coefficient
$1-z_m-a=(m-1)a>0$. The strict inequality holds because the only $m=1$
case is the excluded duplicate. This leaves a three-term sum and gives
the same conclusion.

Finally, if $a<q_m$, the increasing-base signs are $+,-,+,-$, so there
are at most three real zeros. One is $s=0$, and the difference is negative
for large $s$. It is also negative at $s=1,2$ by
\eqref{eq:light-strict-small}. Since its entropy series sums to zero,
it must be positive at some integer $j\ge3$. There is therefore one
crossing between $2$ and $j$, and another after $j$. These, together
with zero, exhaust the permitted zeros and must all be simple. They give
exactly the interval and signs stated in the theorem.
\end{proof}

\subsection{Finite phase intervals}
The crossing results give an alternative exact description of the finite
maximum in Theorem~\ref{thm:finite-classification}. Every endpoint below
is a specified root with a proved crossing orientation.

For a heavy multiplicity $m\ge k$ other than the binary duplicate, put
\begin{equation}\label{eq:phase-cutoff}
 T_m=\tau_{m,m+1},\qquad
 R_m=\max\{m+1,\lceil hT_m+1\rceil\},
\end{equation}
and define the finite extrema
\begin{equation}\label{eq:phase-endpoints}
 \begin{aligned}
 \underline{s}_m
 &=\max\bigl(\{\alpha_m\}\cup
                \{\tau_{j,m}:k\le j<m\}\bigr),\\
 \overline{s}_m
 &=\min\bigl(\{\beta_m\}\cup
                \{\tau_{m,j}:m<j\le R_m\}\bigr).
 \end{aligned}
\end{equation}
The second set contains $\tau_{m,m+1}$, so its minimum is finite even
when $\beta_m=\infty$.

\begin{corollary}[Winning intervals]\label{cor:phase-intervals}
For integer $t\ge1$, the heavy distribution $p_m^{\rm H}$ is globally
optimal if and only if
\begin{equation}\label{eq:heavy-phase}
 \underline{s}_m\le t\le\overline{s}_m.
\end{equation}
If the lower endpoint exceeds the upper endpoint, this multiplicity never
wins. The light or uniform candidate is globally optimal if and only if
\begin{equation}\label{eq:light-phase}
 t\notin
 \bigcup_{\substack{k\le m\le M(t,h)\\p_m^{\rm H}\ne p_{\rm L}}}
       (\alpha_m,\beta_m),
\end{equation}
where equality or inequality of distributions is understood up to
permutation and zero coordinates. All endpoint ties are included.
\end{corollary}
\begin{proof}
Theorem~\ref{thm:heavy-crossing} gives the lower crossing constraints
against smaller heavy multiplicities and the upper constraints against
larger ones. Theorem~\ref{thm:light-crossing} gives the closed interval
$[\alpha_m,\beta_m]$ for comparison with the baseline.

A candidate with multiplicity $m$ loses to $m+1$ when $t>T_m$. For
$1\le t\le T_m$, we have $R_m\ge ht+1$, so
Lemma~\ref{lem:cutoff} implies $b_j(t)<b_{R_m}(t)$ for every $j>R_m$.
It is therefore enough to compare with the finitely many larger
multiplicities in~\eqref{eq:phase-endpoints}. This proves
\eqref{eq:heavy-phase}. The light criterion follows directly from
Theorem~\ref{thm:finite-classification} and the strict sign intervals
in~\eqref{eq:light-crossing-sign}.
\end{proof}

At fixed entropy, heavy winning multiplicities cannot decrease as the
integer sample size increases. Indeed, if $i<j$ and $b_j(t_1)\ge b_i(t_1)$,
then $t_1\ge\tau_{ij}$; for every $t_2>t_1$ we have
$b_j(t_2)>b_i(t_2)$. This assertion concerns comparisons as the sample
size varies. The heavy sequence in $m$ need not be unimodal at a fixed
sample size, as Appendix~\ref{app:certificates} shows.

\section{A fixed finite alphabet}\label{app:finite-alphabet}
Let $B_{t,N}(h)$ denote the maximum when the alphabet has exactly $N$
available symbols; zero coordinates are permitted. If $N=1$ or $h=0$,
only point masses can maximize, and the maximum is zero. Suppose now
that $N\ge2$, $h>0$, and $t\ge1$ is an integer.

Form the saturated candidate list from
Theorem~\ref{thm:finite-classification}, retaining only distributions
with at most $N$ positive coordinates. In addition, let $z$ run through
all roots in $(0,1)$ of
\begin{equation}\label{eq:finite-stationarity}
 f_t'(z)=f_t'\!\left(\frac{1-z}{N-1}\right)
\end{equation}
for which $E_{N-1}(z)\le h$, and add $p_{N-1}(z)$ to the list. Explicitly,
\eqref{eq:finite-stationarity} is the polynomial equation
\begin{equation}\label{eq:finite-polynomial}
 \begin{aligned}
 &(N-1)^t(1-z)^{t-1}\bigl(1-(t+1)z\bigr)\\
 &\qquad-(N-2+z)^{t-1}\bigl((t+1)z+N-t-2\bigr)=0.
 \end{aligned}
\end{equation}
This polynomial is nonzero and has degree at most $t$. The full-support
uniform law is included as the root $z=1/N$ whenever it is feasible.

\begin{theorem}\label{thm:finite-alphabet}
The maximum $B_{t,N}(h)$ is the largest objective value in the finite list
just described. Up to permutation, its maximizing vectors are exactly the
candidates attaining that value.
\end{theorem}
\begin{proof}
Theorems~\ref{thm:main} and~\ref{thm:exceptional} give the possible
shapes. If a maximizer has $H(p)=h$, it is an entropy-matched candidate
whose support fits in $N$. A heavy candidate with multiplicity above
$M(t,h)$ cannot maximize even on this alphabet: the candidate at
$M(t,h)$ has smaller support, the same entropy, and a strictly larger
objective.

If a maximizer has $H(p)<h$, it has full support. Otherwise the
atom-splitting perturbation in the proof of Theorem~\ref{thm:main} is
available on an unused coordinate and strictly improves the objective.
On full support the entropy inequality is locally inactive, so ordinary
stationarity under normalization gives equality of $f_t'$ at every
positive size. This is~\eqref{eq:finite-stationarity} for a nonuniform
one-exceptional vector; a uniform vector satisfies it as well.

Equation~\eqref{eq:finite-polynomial} is obtained by clearing the factor
$(N-1)^t$. It is not identically zero: the derivative difference in
\eqref{eq:finite-stationarity}, extended to $z=1$, is $-2$ for $t=1$
and $-1$ for $t>1$. Thus at most $t$ roots are added. Every listed vector
is feasible, and the list contains every maximizer; attainment proves
the theorem.
\end{proof}

\section{Certified nonmonotonicity examples}\label{app:certificates}
The finite classification need not reduce to a unimodal search in the
heavy multiplicity or to a single transition as entropy increases. The
following strict inequalities are verified by the accompanying
\texttt{missing\_mass\_extremizers\_certificates.py}, which uses only
integer and rational arithmetic from the Python standard library.

At $h=19/8$ and $t=8$, the heavy candidates satisfy
\begin{equation}\label{eq:nonunimodal-example}
 \begin{aligned}
 0.456467562315&<b_{10}(8)<0.456467562316,\\
 0.455820123956&<b_{11}(8)<0.455820123957,\\
 0.457052010607&<b_{12}(8)<0.457052010608.
 \end{aligned}
\end{equation}
In particular $b_{10}(8)>b_{11}(8)<b_{12}(8)$, so the heavy sequence is
not unimodal. The cutoff is $M=20$. Comparing the light candidate and all
heavy candidates $10\le m\le20$ gives the unique winner $m=14$, with
\[
 0.458686581808<B_8(19/8)<0.458686581809.
\]
Its heavy atom lies in
\[
 \frac{343937132264853115880129584166}{10^{30}}
 <z_{14}<
 \frac{343937132264853115880129584167}{10^{30}},
\]
and every other positive atom is $(1-z_{14})/14$.

For $t=3$, the type of the unique maximizer changes from one-light to
one-heavy and back to one-light across the three entropy budgets below.
In the two numerical columns, each displayed decimal $d$ is a strict
lower bound and $d+10^{-12}$ is a strict upper bound.
\begin{center}
\begin{tabular}{c c c c c}
\toprule
$h$ & $\ell(3)$ & $\max_{m\ge k}b_m(3)$ & $M$ & Optimal type\\
\midrule
$11/10$ & $0.296480675541$ & $0.294054230107$ & $5$ & One-light\\
$6/5$   & $0.327755692052$ & $0.333702359978$ & $5$ & One-heavy\\
$7/5$   & $0.424289708030$ & $0.403758744660$ & $6$ & One-light\\
\bottomrule
\end{tabular}
\end{center}
In the middle row the winning heavy multiplicity is $m=3$.

For completeness, the certificate method is as follows. The script stores
a rational bracket $(a,b)$ of width $10^{-30}$ for each entropy root and
verifies the appropriate entropy signs at its endpoints. Logarithms are
bounded after reduction to $x=2^j y$, $1\le y<2$, by setting
$v=(y-1)/(y+1)\in[0,1/3)$ and using
\begin{equation}\label{eq:log-certificate}
 0\le\log y-2\sum_{r=0}^{R-1}\frac{v^{2r+1}}{2r+1}
 \le\frac{2v^{2R+1}}{(2R+1)(1-v^2)}.
\end{equation}
The same formula at $v=1/3$ bounds $\log2$. The script uses $R=40$ and
reverses the appropriate interval endpoints when multiplying by a
negative $j$. For $z\in(a,b)$, positivity of the weights and bases gives
\[
 \begin{aligned}
 V_{m,t}(z)&\ge a(1-b)^t
 +(1-b)\left(1-\frac{1-a}{m}\right)^t,\\
 V_{m,t}(z)&\le b(1-a)^t
 +(1-a)\left(1-\frac{1-b}{m}\right)^t.
 \end{aligned}
\]
These are rational bounds for integer $t$. The stored data certify the
entropy index $k$, the cutoff $M$, every entropy-root bracket, all strict
objective comparisons, and the completeness of each candidate list.

\section*{AI disclosure}
GPT-6 Astra and Claude Opus 5 were used for proof discovery, simplification,
and editorial suggestions. The author is fully responsible for the
correctness of the arguments.

\begin{thebibliography}{2}
\bibitem{BK} D.~Berend and A.~Kontorovich,
\emph{The missing mass problem},
Statistics \& Probability Letters \textbf{82} (2012), no.~6, 1102--1110.
\href{https://doi.org/10.1016/j.spl.2012.02.014}{doi:10.1016/j.spl.2012.02.014}.

\bibitem{BKZ} D.~Berend, A.~Kontorovich, and G.~Zagdanski,
\emph{The expected missing mass under an entropy constraint},
Entropy \textbf{19} (2017), no.~7, article~315, 13~pp.
\href{https://doi.org/10.3390/e19070315}{doi:10.3390/e19070315}.
\end{thebibliography}
\end{document}
